We asked how much machine learning can add to last year's government allocation and what peers and explanations can tell us when there is no last year. On official COFOG data the answer is modest, and more informative than the legacy $R^2>0.99$ suggested: last year's allocation is a very strong baseline, and a learned correction using same-year covariates improves on it by only about \ResPairedGain{} percentage points of the budget for one learner seed (\AnnSeedMeanGain{} pp on average over five seeds, with an interval that excludes zero in \AnnSeedNExcl{} of \AnnSeedN); without history, a stack, a nearest-economies average and a Dirichlet regression beat the income-group mean by about 1.3 pp on the wide ten-division panel but not robustly on the 16-function panel; health and R\&D intensity are predictable in unseen economies, though a ridge regression gets most of the way there; shocks are not anticipated; errors are two to three times larger for lower-income economies; and SHAP and LIME agree on direction more than on ranking and are ambiguous when indicators are correlated. The stacked learner and the bounded-simplex projection are evaluated components, not recommended methods: the stack ties simpler alternatives and loses to persistence in annual mode. The released dataset, protocol, negative results and the audit that explains why an earlier benchmark looked better are, in our view, a more reliable contribution than a higher score.

\textbf{Future work.} (1) Strict forecasting with lagged covariates and publication lags, a rolling-origin evaluation of the headline comparison, and fold-internal imputation. (2) Extended COFOG coverage for non-OECD economies (IMF GFS, national treasuries; India's Union Budget as a case study) and stronger simple cold-start baselines. (3) Re-running Panel B without fiscal-envelope features, and permutation-importance and on-manifold deletion baselines for the explanations. (4) A curated library of statutory floors for the allocation engine. (5) Conformal intervals with group-conditional coverage guarantees for lower-income economies. (6) A human-subject study measuring decision time, agreement and trust calibration of analysts using the explanations.
